% TOP_PROBLEM: {"id":30007676,"problem_number":"LOCAL-30007676","title":"Merger efficiencies","category_id":21,"category":{"id":21,"name":"economics","display_name":"Economics","description":"Open economic theory statements, published research agendas, and proposed scoped specifications; question provenance and historical priority are qualified per record.","slug":"economics","order_index":21,"created_at":"2026-10-08T00:00:00Z"},"status":"research_question","external_url":null,"legacy_ids":[],"published":false,"statement":"Test whether documented merger-efficiency claims predict audited cost savings and consumer price benefits three years later.","economics":{"original_id":165,"field":"Industrial organization and competition","kind":"empirical","formalization_basis":"adapted_research_question","source_status":"research_agenda","priority_status":"earliest_found","priority_note":"Earliest relevant explicit question or agenda passage located in this audit. No claim of first-ever formulation; earlier versions and later solutions were not exhaustively traced.","earliest_verified_reference":{"authors":["Carl Shapiro","Ali Yurukoglu"],"title":"Trends in Competition in the United States: What Does the Evidence Show?","year":2024,"url":"https://web.stanford.edu/~ayurukog/trendsincompetition.pdf","doi":"10.3386/w32762","locator":"Section 6, printed pp. 35–36, merger-efficiency paragraph","evidence_summary":"States room for retrospective studies of whether claimed merger efficiencies materialize."},"current_reference":{"authors":["Carl Shapiro","Ali Yurukoglu"],"title":"Trends in Competition in the United States: What Does the Evidence Show?","year":2024,"url":"https://web.stanford.edu/~ayurukog/trendsincompetition.pdf","doi":"10.3386/w32762","locator":"Section 6, printed pp. 35–36, merger-efficiency paragraph","evidence_summary":"States room for retrospective studies of whether claimed merger efficiencies materialize."},"formalization_note":"The source verifies a broader question or research agenda; the population, estimand, equations and restrictions here are our scoped adaptation. This is not a quoted canonical conjecture, and the audit does not prove contemporary unsolved status for every special case."},"statusQualification":"Published question or research agenda verified at the cited locator. The mathematical specification is adapted for this catalogue and includes additional assumptions; source publication does not establish first-ever priority or certify this exact model remains unsolved. The source verifies a broader question or research agenda; the population, estimand, equations and restrictions here are our scoped adaptation. This is not a quoted canonical conjecture, and the audit does not prove contemporary unsolved status for every special case.","statusReviewedAt":"2026-10-08"}
% FIRST_POSTED: 2026-10-08
% ENTRY_KIND: statement_only
% !TeX program = lualatex
\documentclass[11pt]{article}
\usepackage{fontspec}
\usepackage[a4paper,margin=25mm]{geometry}
\usepackage{amsmath,amssymb,mathtools}
\usepackage{xurl}
\usepackage[hidelinks]{hyperref}
\newcommand{\sref}[2]{\hyperlink{source:#1:#2}{[S#2]}}
\setlength{\emergencystretch}{3em}
\begin{document}
% problemId:30007676
\section*{Merger efficiencies}\label{30007676}
\subsection{Definitions and mathematical statement}
Consider proposed horizontal manufacturing mergers reviewed in the United States over a prespecified five-year cohort. For deal \(m\), let \(S_m\) be claimed annual marginal-cost savings documented before approval, and \(X_m\) other premerger information. Let \(C_m(a,h)\) denote the cost of producing a fixed reference output vector at horizon \(h=3\) years under merger regime \(a\in\{0,1\}\), and let \(P_m(a,h)\) be the corresponding expenditure-weighted consumer price index after equilibrium repricing. Define realized cost savings \(G_m=C_m(0,h)-C_m(1,h)\), the claim-calibration function \(g(s,x)=E[G_m\mid S_m=s,X_m=x]\), and consumer price benefit \(B_m=P_m(0,h)-P_m(1,h)\). Costs and prices are measured in constant dollars; expectation refers to the reviewed-deal population. Obtain production and logistics audits that distinguish technological savings from purchasing-power transfers and output changes. Use approved and appropriately comparable abandoned or blocked deals to construct counterfactuals, recording the selection assumptions needed for each comparison. Estimate the joint distribution of \(G_m\) and \(B_m\), rather than attributing all profit growth to efficiency. Determine which preapproval evidence predicts positive, merger-specific savings and consumer benefits. Where approval selection precludes identification, provide bounds and specify the additional information that would reduce them.

\par\textbf{Formulation and scope.} The source verifies a broader question or research agenda; the population, estimand, equations and restrictions here are our scoped adaptation. This is not a quoted canonical conjecture, and the audit does not prove contemporary unsolved status for every special case.

\par\textbf{Open-question provenance.} An explicit question or research agenda was verified at the source locator. This does not by itself certify that every later version remains unresolved \sref{30007676}{1}.

\par\textbf{Historical priority.} Earliest relevant explicit question or agenda passage located in this audit. No claim of first-ever formulation; earlier versions and later solutions were not exhaustively traced.

\subsection{Short English statement}
Test whether documented merger-efficiency claims predict audited cost savings and consumer price benefits three years later.

\subsection{Sources}
\begin{itemize}\raggedright
\item[\textnormal{[S1]}]\hypertarget{source:30007676:1}{} Carl Shapiro, Ali Yurukoglu. 2024. Trends in Competition in the United States: What Does the Evidence Show?. Section 6, printed pp. 35–36, merger-efficiency paragraph.
\url{https://web.stanford.edu/~ayurukog/trendsincompetition.pdf}.
DOI: 10.3386/w32762.
{\small\textit{Earliest verified question/agenda reference; historical priority qualified below. States room for retrospective studies of whether claimed merger efficiencies materialize.}}
\end{itemize}
\end{document}
